\section{Method}
\label{sec:method}

% TODO: Add architecture overview figure
% \begin{figure*}[t]
%   \centering
%   \includegraphics[width=\linewidth]{fig/overview.pdf}
%   \caption{Overview of our approach. (a) Each modality is processed by SAM2's frozen image encoder augmented with Soft-MoE LoRA layers. (b) A cross-modal fusion head jointly compares all modality features and produces softmax-normalized relative weights. (c) These weights are max-normalized for memory modulation (UAMM) and used as raw softmax for output fusion (AMF).}
%   \label{fig:overview}
% \end{figure*}

Our method adapts MemorySAM~\cite{chen2025memorysam} for the MULTIAQUA maritime segmentation task.
We retain the core idea of treating each modality as a sequential frame processed through SAM2's~\cite{ravi2024sam2} memory-based tracking pipeline, and modify the LoRA adaptation strategy, modality scoring, and output fusion.
Our pipeline operates in four phases: given $M$ input modalities (RGB, thermal, LiDAR), (1) all modalities are first encoded by SAM2's frozen image encoder with Soft-MoE LoRA layers; (2) a cross-modal fusion head jointly compares the backbone features of all modalities and produces softmax-normalized relative weights; (3) these weights are max-normalized and used to modulate vision features before SAM2's memory-augmented tracking step (UAMM); and (4) the per-modality mask outputs are fused using the raw softmax weights (AMF).


\subsection{Soft-MoE LoRA}
\label{sec:soft_moe_lora}

We keep SAM2's image encoder frozen and inject low-rank adaptation matrices into the query and value projections of each attention block, following the standard LoRA~\cite{hu2022lora} formulation.
Rather than using a single pair of low-rank matrices, we maintain $N$ expert pairs $\{(\mathbf{A}_i, \mathbf{B}_i)\}_{i=1}^{N}$, where $\mathbf{A}_i \in \mathbb{R}^{d \times r}$ and $\mathbf{B}_i \in \mathbb{R}^{r \times d}$, with $r$ being the LoRA rank and $d$ the feature dimension.

Following the Soft Mixture-of-Experts formulation~\cite{puigcerver2024softmoe}, instead of selecting the top-$k$ experts via hard routing~\cite{zhu2024moe, shazeer2017moe}, we compute a softmax-weighted combination over all experts for each input token $\mathbf{x}$:
\begin{equation}
  \Delta \mathbf{x} = \sum_{i=1}^{N} \alpha_i(\mathbf{x}) \cdot \mathbf{B}_i \mathbf{A}_i \mathbf{x},
  \quad \alpha_i(\mathbf{x}) = \frac{\exp(g_i(\mathbf{x}))}{\sum_{j=1}^{N} \exp(g_j(\mathbf{x}))},
  \label{eq:soft_moe}
\end{equation}
where $g(\mathbf{x}) = \mathbf{W}_g \mathbf{x} + \mathbf{b}_g$ is a lightweight gating network with $\mathbf{W}_g \in \mathbb{R}^{N \times d}$.
The output of the original QKV projection is augmented as $\mathbf{q} \leftarrow \mathbf{q} + \Delta \mathbf{q}$ and $\mathbf{v} \leftarrow \mathbf{v} + \Delta \mathbf{v}$.

When the number of experts $N$ is small (\eg, $N{=}3$ matching three modalities), hard top-$k$ routing with $k{=}2$ permanently disables one expert per token.
Soft gating avoids this issue by ensuring gradient flow to every expert at every step, enabling all experts to specialize without the dead expert problem.
We use $N{=}3$ experts and $r{=}4$ throughout our experiments.


\subsection{Quality-Aware Cross-Modal Scoring}
\label{sec:cross_modal}

A central challenge in multi-modal fusion is estimating the relative quality of each modality.
Prior approaches~\cite{chen2025memorysam} use independent per-modality heads (\eg, a CNN followed by sigmoid activation) to produce scalar confidence scores.
However, when the backbone features are rich and all modalities produce large logits, the sigmoid function saturates---all scores converge to ${\sim}1.0$ regardless of actual quality differences.
This effectively disables modality-adaptive weighting.
Moreover, global average pooling alone encodes semantic content rather than quality: a degraded RGB frame and a clear RGB frame may yield similar GAP vectors after backbone processing, preventing the network from identifying which modality is degraded.

We address this by enriching each modality's representation with quality-sensitive statistics.
Given backbone feature maps $\{\mathbf{F}_i\}_{i=1}^{M}$ from Phase~1, we compute a multi-pool descriptor for each modality:
\begin{equation}
  \mathbf{q}_i = [\operatorname{GAP}(\mathbf{F}_i)\,;\,\operatorname{GMP}(\mathbf{F}_i)\,;\,\sigma(\mathbf{F}_i)],
  \label{eq:multipool}
\end{equation}
where $\operatorname{GAP}$ and $\operatorname{GMP}$ denote global average and max pooling respectively, and $\sigma(\mathbf{F}_i) \in \mathbb{R}^{C}$ is the channel-wise standard deviation computed over spatial locations.
Together, the triple $(\operatorname{GAP}, \operatorname{GMP}, \sigma)$ captures mean activation level, peak response, and spatial variability---statistics that change when a sensor degrades.

Each modality's descriptor $\mathbf{q}_i \in \mathbb{R}^{3C}$ is compressed by a modality-specific network $\phi_i$:
\begin{equation}
  \mathbf{z}_i = \phi_i(\mathbf{q}_i) = \operatorname{ReLU}(\mathbf{W}_c^{(i)} \cdot \mathbf{q}_i + \mathbf{b}_c^{(i)}),
  \label{eq:compress}
\end{equation}
where $\mathbf{W}_c^{(i)} \in \mathbb{R}^{h \times 3C}$ are modality-specific weights, allowing each sensor type to learn its own quality-relevant feature extraction.
The compressed vectors are concatenated and passed through a comparison MLP to produce modality weights:
\begin{equation}
  \mathbf{l} = \psi([\mathbf{z}_1; \mathbf{z}_2; \ldots; \mathbf{z}_M]),
  \quad w_i = \frac{\exp(l_i / \tau)}{\sum_{j=1}^{M} \exp(l_j / \tau)},
  \label{eq:cross_modal}
\end{equation}
where $\psi$ is a two-layer MLP and $\tau$ is a temperature parameter.
The softmax normalization inherently produces relative weights: even when all modalities are high quality, the network can express fine-grained preferences; when one modality degrades, its weight decreases and others increase.
We initialize the output layer of $\psi$ to zero, so that $\mathbf{l} = \mathbf{0}$ at the start of training and all modalities receive uniform weight $1/M$, ensuring stable initialization.


\subsection{Gating Oracle Supervision}
\label{sec:oracle}

The segmentation objective alone provides an insufficient gradient signal to drive adaptive gating.
A network that assigns fixed average weights (\eg, $\mathbf{w} \approx [0.27, 0.36, 0.37]$) already minimizes the main segmentation loss because the downstream mask fusion can compensate.
We observed this in practice: without additional supervision, the learned weights converge to near-constant values regardless of input, effectively behaving as a static weighted average.

To address this, we attach a lightweight auxiliary segmentation head $\mathcal{H}_i$ to each modality's backbone features.
Each head independently predicts the full segmentation map for its modality, providing a per-image estimate of that modality's discriminative quality.
We derive oracle gating weights from these auxiliary predictions:
\begin{equation}
  \tilde{w}_i = \frac{\exp(-\mathcal{L}_\text{aux}^{(i)})}{\sum_{j=1}^{M} \exp(-\mathcal{L}_\text{aux}^{(j)})},
  \label{eq:oracle}
\end{equation}
where $\mathcal{L}_\text{aux}^{(i)}$ is the per-image cross-entropy loss for modality $i$'s auxiliary head.
Intuitively, the oracle assigns higher weight to modalities whose backbone features independently support accurate segmentation, and lower weight when a modality is degraded (\eg, dark RGB at night yields high $\mathcal{L}_\text{aux}^\text{rgb}$, reducing $\tilde{w}_\text{rgb}$).

The fusion head is trained to match these oracle weights via KL divergence:
\begin{equation}
  \mathcal{L}_\text{gate} = D_\text{KL}(\tilde{\mathbf{w}} \,\|\, \mathbf{w}),
  \label{eq:gate_kl}
\end{equation}
where $\tilde{\mathbf{w}}$ is treated as a fixed target (detached from the computation graph).
This provides a direct, per-image training signal that forces the fusion head to respond to input quality variation rather than converging to static modality preferences.


\subsection{Uncertainty-Aware Memory Modulation}
\label{sec:uamm}

Before each modality's features enter SAM2's memory-augmented tracking step, we modulate them based on the cross-modal weights.
To preserve the best modality's features at full strength while suppressing degraded ones, we apply max-normalization to the softmax weights:
\begin{equation}
  \hat{w}_i = \frac{w_i}{\max_{j} w_j + \epsilon},
  \quad \hat{\mathbf{f}}_i^{(k)} = \hat{w}_i \cdot \mathbf{f}_i^{(k)}, \quad \forall k,
  \label{eq:uamm}
\end{equation}
where $\mathbf{f}_i^{(k)}$ are the multi-scale vision features for modality $i$ and $\epsilon = 10^{-8}$.

The max-normalization ensures that the highest-scoring modality always receives $\hat{w} = 1.0$ (no attenuation), while other modalities are scaled down proportionally.
This addresses the sigmoid saturation problem: unlike sigmoid-based gating where all scores can converge to ${\sim}1.0$, the softmax-based relative comparison always differentiates between modalities, and max-normalization preserves the absolute feature magnitude of the best modality.
For instance, if thermal is the most reliable modality at night ($w_\text{thermal} = 0.6$, $w_\text{lidar} = 0.3$, $w_\text{rgb} = 0.1$), the UAMM scores become $\hat{w}_\text{thermal} = 1.0$, $\hat{w}_\text{lidar} = 0.5$, $\hat{w}_\text{rgb} = 0.17$---thermal features pass through unmodified while RGB is strongly suppressed.

The modulated features $\hat{\mathbf{f}}_i$ are then processed by SAM2's \texttt{track\_step}, which updates the memory bank with cross-attention.
The first modality serves as the initialization frame, and subsequent modalities attend to accumulated memory, enabling cross-modal feature exchange through SAM2's built-in temporal attention mechanism.


\subsection{Adaptive Mask Fusion}
\label{sec:amf}

After obtaining per-modality mask predictions $\mathbf{o}_i$ from the tracking loop, we fuse them using the raw softmax weights from the cross-modal fusion head:
\begin{equation}
  \mathbf{o}_\text{fused} = \sum_{i=1}^{M} w_i \cdot \mathbf{o}_i,
  \label{eq:amf}
\end{equation}
where $w_i$ are the softmax weights from \cref{eq:cross_modal}, which already sum to one.

This dual-path design---max-normalized weights for UAMM, raw softmax weights for AMF---serves complementary purposes.
UAMM requires \textit{absolute} preservation: the best modality should contribute its full feature magnitude to the memory bank.
AMF requires \textit{relative} weighting: each modality's mask output should be weighted proportionally to its estimated quality.
Reusing the same cross-modal fusion head for both paths avoids redundant computation while tailoring the weight application to each stage's requirements.


\subsection{Prototype Segmentation Loss}
\label{sec:proto_loss}

Following MemorySAM~\cite{chen2025memorysam}, we use a Semantic Prototype Memory Module to guide SAM2 toward semantic-level understanding.
For each class $c$, a global prototype $\mathbf{p}_c$ is maintained via exponential moving average (momentum $= 0.8$).
At each training step, batch prototypes $\hat{\mathbf{p}}_c$ are computed by averaging the fused feature vectors $\mathbf{F}_\text{fused}$ belonging to class $c$ (determined by resizing the ground-truth labels to the feature map resolution).
The prototype loss is defined as:
\begin{equation}
  \mathcal{L}_\text{proto} = \text{MSE}(\hat{\mathbf{p}}, \mathbf{p}),
  \label{eq:proto}
\end{equation}
and the total training loss is:
\begin{equation}
  \mathcal{L} = \mathcal{L}_\text{OHEM} + \lambda_\text{proto} \cdot \mathcal{L}_\text{proto} + \lambda_\text{gate} \!\left( \mathcal{L}_\text{gate} + \gamma \cdot \frac{1}{M}\sum_{i=1}^{M} \mathcal{L}_\text{aux}^{(i)} \right),
  \label{eq:total_loss}
\end{equation}
where $\mathcal{L}_\text{OHEM}$ is the Online Hard Example Mining cross-entropy loss, $\gamma{=}0.3$ controls the auxiliary segmentation contribution, and $\lambda_\text{gate}{=}0.5$ in all experiments.


\subsection{Night-Robust RGB Augmentation}
\label{sec:night_aug}

The MULTIAQUA dataset requires generalization from daytime training to nighttime evaluation, where RGB quality degrades severely while thermal and LiDAR remain relatively stable.
Inspired by the augmentation strategies proposed in the MULTIAQUA paper~\cite{muhovic2025multiaqua} and CRM~\cite{shin2024crm}, we apply three complementary RGB corruptions during training (applied only to the RGB modality, leaving thermal and LiDAR unchanged):

\noindent\textbf{Night Simulation} (probability 0.35): extreme brightness reduction (factor sampled from $[0.03, 0.25]$), followed by contrast reduction ($[0.3, 0.7]$), gamma correction ($\gamma \in [0.4, 0.8]$), and additive Gaussian noise ($\sigma = 0.02$).

\noindent\textbf{Complementary Random Masking} (probability 0.30): rectangular patches are randomly placed on the RGB image, zeroing out 20--50\% of spatial regions.
This forces the model to rely on thermal and LiDAR for the masked areas.

\noindent\textbf{Full-Channel Zeroing} (probability 0.12): the entire RGB input is replaced with zeros, simulating complete RGB failure.
This is the most extreme case, requiring the model to produce reasonable segmentation from thermal and LiDAR alone.

These augmentations are stacked sequentially: an image may undergo night simulation and then additionally be partially masked, creating diverse degradation patterns.
